\item {\bf Product Ethics Exploration: Benchmarking}

You will choose a modern AI product to analyze throughout the assignments and gain initial access to it. Please note that the questions regarding your chosen product are interconnected and will progressively build upon each other from one assignment to the next.

\textbf{Suggested resources:}
\begin{itemize}
    \item Pick an AI-powered product or tool (not just a raw LLM chatbot) that is accessible to you. Good candidates include: Cursor, Nano Banana, ElevenLabs, Perplexity, Claude Code, etc. Free trials or student-accessible tools are recommended.
    \item The product's official website and onboarding/getting started pages
    \item Privacy policy, Terms of Service, Acceptable Use Policy
    \item App store listings, GitHub documentation or READMEs (if open-source or developer-targeted)
    \item Recent news coverage (e.g., TechCrunch, The Verge, Wired, NYT technology section)
    \item Reddit communities, Hacker News, or user forums for first impressions and experiences
\end{itemize}

\textbf{In this part:} You will analyze the company's policies around data collection, transparency, copyright, and user rights.

\vspace{2mm}

\textbf{Why it's important:} These policies determine how user data is used, what transparency exists around the AI system, and what legal and ethical frameworks govern the product's operation.

\vspace{2mm}

\textbf{Suggested resources:}
\begin{itemize}
  \item Terms of Service and Privacy Policy
  \item Data usage pages (e.g., ``How we use your data'')
  \item Security or compliance documentation (GDPR, CCPA statements)
  \item Company blog posts about transparency, trust, and safety
  \item Public reports of harms or incidents (news articles, watchdog groups)
\end{itemize}

\begin{enumerate}
    \input{06-product-ethics/01-data-collection}
    \input{06-product-ethics/02-trust}
    \input{06-product-ethics/03-fairness}
    \input{06-product-ethics/04-accountability}
    \input{06-product-ethics/05-copyright}
\end{enumerate}